% Irodov Part III §3.1 original 41--53, PDF pp.106--107,
% printed pp.97--98; source figures 3.5--3.7 checked.
\electroproblem{41}{97}
点偶极矩 $\bm p$ 与均匀外电场 $\bm E_0$
方向相同。证明总电场有一张
以偶极子为球心的球形等势面，
求该球半径。
\exsolution
取 $z$ 轴沿 $\bm p$ 与 $\bm E_0$ 方向。
点偶极子电势
$\varphi_d=kp\cos\theta/r^2$，
均匀电场电势（选零点在原点）
$\varphi_0=-E_0r\cos\theta$。
总电势
\[
\varphi(r,\theta)=
\cos\theta\left(\frac{kp}{r^2}-E_0r\right).
\]
若括号为零，对\emph{全部}极角 $\theta$
有 $\varphi=0$，形成球形等势面。
由 $kp/R^2=E_0R$，
\[
\boxed{R=\left(\frac{p}{4\pi\varepsilon_0E_0}\right)^{1/3}}.
\]
\excheck 半径需取\emph{立方根}；
PDF 文字抽取容易把根号前的3写成乘数。

\electroproblem{42}{97}
两根相互平行、间距 $\ell$ 的长直细线，
线电荷密度分别为 $\lambda,-\lambda$。
求远离两线、距中点 $r\gg\ell$
且与连线夹角为 $\theta$ 处的电势和电场模。
\exsolution
在垂直两线的截面上，选两线中点为原点，
正负线位于 $\pm\bm \ell/2$。
二维线电荷的电势差叠加为
\[
\varphi=\frac{\lambda}{2\pi\varepsilon_0}
\ln\frac{r_-}{r_+}.
\]
在 $r\gg\ell$ 时，
$\ln(r_-/r_+)\simeq(\ell/r)\cos\theta$，
因此
\[
\boxed{\varphi(r,\theta)\simeq
\frac{\lambda\ell\cos\theta}{2\pi\varepsilon_0r}}.
\]
其极坐标场分量为
\[
E_r=-\partial_r\varphi=
\frac{\lambda\ell\cos\theta}{2\pi\varepsilon_0r^2},
\quad
E_\theta=-r^{-1}\partial_\theta\varphi=
\frac{\lambda\ell\sin\theta}{2\pi\varepsilon_0r^2}.
\]
故
\[
\boxed{E\simeq\frac{|\lambda|\ell}{2\pi\varepsilon_0r^2}}.
\]
\excheck 势依赖方向 $\cos\theta$，
但此二维偶极近场阶的\emph{场强模}
不依赖角度。

\electroproblem{43}{97}
两个同轴均匀带电细圆环，半径均为 $R$，
总电荷分别为 $+q,-q$，
中心相距 $\ell\ll R$。
按图3.6取中点为轴坐标 $x=0$，
求轴上电势 $\varphi(x)$ 和 $E_x(x)$
及其大致图形，并讨论 $|x|\gg R$。
\exsolution
令正环位于 $x=+\ell/2$，
负环位于 $x=-\ell/2$。
轴上电势精确表达式
\[
\varphi(x)=kq\left[
\frac1{\sqrt{R^2+(x-\ell/2)^2}}
-\frac1{\sqrt{R^2+(x+\ell/2)^2}}\right].
\]
对 $\ell/R\ll1$ 作一阶泰勒展开：
\[
\boxed{\varphi(x)\simeq
\frac{q\ell\,x}{4\pi\varepsilon_0
(R^2+x^2)^{3/2}}}.
\]
由 $E_x=-\mathrm d\varphi/\mathrm dx$ 得
\[
\boxed{E_x(x)\simeq
\frac{q\ell(2x^2-R^2)}
{4\pi\varepsilon_0(R^2+x^2)^{5/2}}}.
\]
$\varphi(x)$ 是奇函数，在 $x=0$ 过零，
在 $x=\pm R/\sqrt2$ 处有极值；
$E_x(x)$ 是偶函数，
中心为负，$|x|>R/\sqrt2$ 为正，
且远方趋于零。
当 $|x|\gg R$，
\[
\boxed{\varphi\sim\frac{q\ell\,\operatorname{sgn}x}
{4\pi\varepsilon_0x^2}},\qquad
\boxed{E_x\sim\frac{q\ell}{2\pi\varepsilon_0|x|^3}}.
\]
\begin{center}
\begin{tikzpicture}[scale=0.56,>=Stealth]
\draw[->](-3.4,0)--(3.4,0) node[right]{$x/R$};
\draw[->](0,-1.35)--(0,1.35) node[above]{$\varphi$};
\draw[thick,domain=-3:3,samples=60]
plot(\x,{2.6*\x/((1+\x*\x)^1.5)});
\begin{scope}[shift={(8,0)}]
\draw[->](-3.4,0)--(3.4,0) node[right]{$x/R$};
\draw[->](0,-1.35)--(0,1.35) node[above]{$E_x$};
\draw[thick,domain=-3:3,samples=60]
plot(\x,{1.8*(2*\x*\x-1)/((1+\x*\x)^2.5)});
\end{scope}
\end{tikzpicture}
\end{center}
\excheck 图线为无量纲形状示意。
远场是电偶极矩 $p=q\ell$ 沿环轴方向，
因此 $\varphi\propto |x|^{-2}$、
$E\propto|x|^{-3}$。

\electroproblem{44}{97}
两张无限大平面面电荷密度为
$+\sigma,-\sigma$，间距 $\ell\ll R$，
各有半径 $R$ 的同轴圆孔（图3.7）。
以孔轴中点为原点，沿轴取 $x$，
求电势 $\varphi(x)$ 与场分量 $E_x(x)$，
说明电势图像。
\exsolution
将“带圆孔平面”分解为
无限均匀带电平面，减去同号的带电圆盘。
两平面分离 $\ell\ll R$，
用薄层偶极展开并选
$\varphi(0)=0$，得到
\[
\boxed{\varphi(x)\simeq
\frac{\sigma\ell}{2\varepsilon_0}
\frac{x}{\sqrt{x^2+R^2}}}.
\]
对其求导：
\[
\boxed{E_x(x)=-\frac{\mathrm d\varphi}{\mathrm dx}
=-\frac{\sigma\ell R^2}
{2\varepsilon_0(x^2+R^2)^{3/2}}}.
\]
电势为单调增加的奇函数，
在 $x\to\pm\infty$ 时趋于
$\pm\sigma\ell/(2\varepsilon_0)$，
电场为关于原点对称的负向单峰。
\begin{center}
\begin{tikzpicture}[scale=0.67,>=Stealth]
\draw[->](-3.7,0)--(3.7,0) node[right]{$x/R$};
\draw[->](0,-1.3)--(0,1.3) node[above]{$\varphi$};
\draw[dashed](-3.3,1)--(3.3,1);
\draw[dashed](-3.3,-1)--(3.3,-1);
\draw[thick,domain=-3.3:3.3,samples=50]
plot(\x,{\x/sqrt(1+\x*\x)});
\end{tikzpicture}
\end{center}
\excheck 这里选择了图3.7中的面电荷符号；
若将正负平面互换，
$\varphi,E_x$ 整体反号，但模和函数形状不变。

\electroproblem{45}{98}
两张半径 $R$ 的平行圆盘相距
$\ell\ll R$，
面电荷密度分别为 $+\sigma,-\sigma$。
求沿轴线、离两圆盘距离
$x\gg\ell$ 处的电势和电场模，
再讨论 $x\gg R$。
\exsolution
一张均匀带电圆盘在轴上的电势为
$\varphi_{\rm disk}(z)=
[\sigma/(2\varepsilon_0)]
(\sqrt{R^2+z^2}-|z|)$。
分别对两个盘取正负号并用
$\ell\ll x$ 展开其轴向位移，
得到正盘外侧
\[
\boxed{\varphi_+(x)\simeq
\frac{\sigma\ell}{2\varepsilon_0}
\left(1-\frac{x}{\sqrt{x^2+R^2}}\right)},
\]
负盘外侧的电势取相反符号。
沿轴场强模由电势微分给出：
\[
\boxed{E(x)\simeq
\frac{|\sigma|\ell R^2}
{2\varepsilon_0(x^2+R^2)^{3/2}}}.
\]
远处 $x\gg R$，定义电偶极矩
$p=(\pi R^2\sigma)\ell$，于是
\[
\boxed{|\varphi|\simeq
\frac{|p|}{4\pi\varepsilon_0x^2}},\qquad
\boxed{E\simeq\frac{2|p|}{4\pi\varepsilon_0x^3}}.
\]
\excheck 圆盘尺寸仍远小于观察距离时才可
将有限平行板电容器视为点电偶极子。

\electroproblem{46}{98}
点偶极矩 $\bm p$ 位于距无限长均匀线电荷
$r$ 处，线电荷密度为 $\lambda$。
求作用力，当 $\bm p$ 分别：
(a) 沿直线；(b) 沿径向；
(c) 垂直于直线和径向。
\exsolution
线电荷外场为
$\bm E(\bm r_\perp)=
[\lambda/(2\pi\varepsilon_0)]\bm r_\perp/r^2$。
固定偶极矩的受力
$\bm F=(\bm p\cdot\nabla)\bm E$。
(a) 若沿线方向 $\bm p=p\bm e_z$，
场与 $z$ 无关，
\[
\boxed{\bm F_a=0}.
\]
(b) 若 $\bm p=p\bm e_r$，
径向求导
$\mathrm d(\lambda/(2\pi\varepsilon_0r))/\mathrm dr
=-\lambda/(2\pi\varepsilon_0r^2)$，
\[
\boxed{\bm F_b=
-\frac{\lambda p}{2\pi\varepsilon_0r^2}\bm e_r}.
\]
(c) 若 $\bm p=p\bm e_\varphi$，
在固定观察位置处沿切向微分
$\bm E$, 结果沿切向同向：
\[
\boxed{\bm F_c=
\frac{\lambda p}{2\pi\varepsilon_0r^2}\bm e_\varphi}.
\]
\excheck 受力是场梯度作用，
不是均匀场中的 $qE$；
若改变偶极矩的指向，受力也相应反向。

\electroproblem{47}{98}
两个水分子相距
$\ell=10\,\mathrm{nm}$，
各自电偶极矩大小
$p=0.62\times10^{-29}\,\mathrm{C\,m}$，
并沿两分子连线同向排列。
求两分子间相互作用力。
\exsolution
同轴同向的两点偶极子的相互作用势能为
\[
U(\ell)=
\frac{1}{4\pi\varepsilon_0\ell^3}
[\bm p_1\cdot\bm p_2
-3(\bm p_1\cdot\bm e_r)(\bm p_2\cdot\bm e_r)]
=-\frac{2p^2}{4\pi\varepsilon_0\ell^3}.
\]
相对连线方向的力
$F_r=-\mathrm dU/\mathrm d\ell$，
\[
F_r=-\frac{6p^2}{4\pi\varepsilon_0\ell^4}.
\]
所以力为吸引力，大小
\[
\boxed{|F|=
\frac{3p^2}{2\pi\varepsilon_0\ell^4}
\simeq2.1\times10^{-16}\,\mathrm N}.
\]
\excheck 假定两水分子可视为定向的点偶极子，
不是在自由转动分子的热平均取向下计算。

\electroproblem{48}{98}
电场给定为
$\bm E=a(y\bm i+x\bm j)$，
求电势 $\varphi(x,y)$。
\exsolution
由 $E_x=-\partial_x\varphi=-ay$，
对 $x$ 积分
$\varphi=-axy+f(y)$。
再用 $E_y=-\partial_y\varphi=-ax$，
得到 $f'(y)=0$，
因此
\[
\boxed{\varphi(x,y)=-axy+C}.
\]
\excheck $\partial_yE_x=\partial_xE_y=a$，
满足静电场无旋性，电势存在。

\electroproblem{49}{98}
电场
$\bm E=2axy\,\bm i+a(x^2-y^2)\bm j$，
求电势 $\varphi(x,y)$。
\exsolution
由
$\partial_x\varphi=-2axy$，
积分得
$\varphi=-ax^2y+f(y)$。
第二分量要求
$-\partial_y\varphi=a(x^2-y^2)$，
代入
$ax^2-f'(y)=ax^2-ay^2$，
所以 $f'(y)=ay^2$、
$f(y)=ay^3/3+C$。
最终
\[
\boxed{\varphi(x,y)=
ay\left(\frac{y^2}{3}-x^2\right)+C}.
\]
\excheck 两个场分量都由该电势的负梯度恢复，
可用 $\partial_yE_x=\partial_xE_y=2ax$
作独立检查。

\electroproblem{50}{98}
电场
\[
\bm E=ay\,\bm i+(ax+bz)\bm j+by\,\bm k,
\]
其中 $a,b$ 常数。求三维电势。
\exsolution
由 $\partial_x\varphi=-ay$，
先积分
$\varphi=-axy+f(y,z)$。
由 $E_y=-(\partial_y\varphi)
=ax+bz$，
得
$\partial_yf=-bz$，
所以
$f=-byz+g(z)$。
最后由
$E_z=-\partial_z\varphi=by-g'(z)$
得 $g'(z)=0$，
故
\[
\boxed{\varphi(x,y,z)=-y(ax+bz)+C}.
\]
\excheck 逐项负梯度给出原题全部三个分量，
表明输入场确为保守静电场。

\electroproblem{51}{98}
某区域电势只与 $x$ 有关：
$\varphi(x)=-ax^3+b$。
求该区域体电荷密度 $\rho(x)$。
\exsolution
泊松方程为
$\nabla^2\varphi=-\rho/\varepsilon_0$。
电势只依赖 $x$，所以
\[
\frac{\mathrm d^2\varphi}{\mathrm dx^2}
=\frac{\mathrm d}{\mathrm dx}(-3ax^2)=-6ax.
\]
由泊松方程
\[
\boxed{\rho(x)=-\varepsilon_0\nabla^2\varphi
=6\varepsilon_0ax}.
\]
\excheck 同时
$E_x=-\varphi_x=3ax^2$，
$\varepsilon_0\,\mathrm dE_x/\mathrm dx
=6\varepsilon_0ax$，与高斯定律一致。

\electroproblem{52}{98}
两块很大平行平板相距 $d$，
板间空间均匀填充体电荷密度 $\rho$，
板间电势差大小为 $|\Delta\varphi|$。
求使一块板附近电场强度为零的 $\rho$，
及另一板附近的场强大小。
\exsolution
沿板法线取 $x\in[0,d]$，
泊松方程等价于
$\mathrm dE_x/\mathrm dx=\rho/\varepsilon_0$。
若在 $x=0$ 处 $E_x(0)=0$，则
\[
E_x(x)=\frac{\rho x}{\varepsilon_0}.
\]
电势差
$\varphi(d)-\varphi(0)
=-\int_0^dE_x(x)\,\mathrm dx
=-\frac{\rho d^2}{2\varepsilon_0}$。
所以密度的\emph{大小}
\[
\boxed{|\rho|=
\frac{2\varepsilon_0|\Delta\varphi|}{d^2}},
\]
另一板邻近场强模为
\[
\boxed{|E(d)|=\frac{|\rho|d}{\varepsilon_0}
=\frac{2|\Delta\varphi|}{d}}.
\]
\excheck $\rho$ 的代数符号由
$\Delta\varphi$ 的端点定义确定；
本题若未指定电势差方向，
只能无歧义地报告其大小。

\electroproblem{53}{98}
带电实心球内部电势为
$\varphi(r)=ar^2+b$，
$a,b$ 为常数。
求球内空间电荷密度 $\rho(r)$。
\exsolution
径向球对称电势的拉普拉斯算符为
\[
\nabla^2\varphi=
\frac1{r^2}\frac{\mathrm d}{\mathrm dr}
\left(r^2\frac{\mathrm d\varphi}{\mathrm dr}\right).
\]
因为 $\varphi'(r)=2ar$，
\[
\nabla^2\varphi
=\frac1{r^2}\frac{\mathrm d(2ar^3)}{\mathrm dr}
=6a.
\]
由泊松方程
$\nabla^2\varphi=-\rho/\varepsilon_0$，
得到
\[
\boxed{\rho(r)=-6a\varepsilon_0=\text{常数}}.
\]
\excheck 常数电势偏移 $b$ 不产生电场或电荷；
体密度均匀，与球心正则性一致。
